\documentclass[12pt]{article}
\usepackage[utf8]{inputenc}
\usepackage[T1]{fontenc}
\usepackage{hyperref}
\title{Yolngu Sign Language Language Resources for K-12 Curriculum Development}
\author{Compiled by Claude Sonnet 5.5 (claude-sonnet-5-5) for Andr\'as Kornai}
\date{}
\begin{document}
\maketitle

\section{Introduction}
Yolngu Sign Language (YSL, ISO 639-3: ygs) is a shared sign language of north-east Arnhem Land, Australia. Signs accompany spoken Yolngu Matha and are used in ceremony, dance, hunting and situations that call for quiet. Search results describe it as endangered and, until recently, little documented. The search found no school curriculum in YSL. The main documentation is a handbook project.

\section{Primary Educational Resources}
\begin{itemize}
\item \emph{The Illustrated Handbook of Yolngu Sign Language of North East Arnhem Land}. It was developed over about 25 years and directed by Yolngu signers. Descriptions say it presents more than 500 signs in image sequences, with a grammar section, a learner's guide and an English index. The page found is a University of Sydney talk page about the handbook, not the book itself: \url{https://www.sydney.edu.au/museum/whats-on/talks-and-events/previous-talks-and-events/the-illustrated-handbook-of-yolu-sign-language.html}
\item Australia Awards in Indonesia page about the handbook: \url{https://www.australiaawardsindonesia.org/content/630/15/the-illustrated-handbook-of-yolnu-sign-language?sub=true}
\end{itemize}

\section{Linguistic and Reference Materials}
\begin{itemize}
\item Repository record that the search returned for YSL research. It is associated with the Maypilama and Adone 2013 paper, which describes YSL as undocumented. The record's content was not verified: \url{https://clok.uclan.ac.uk/14387}
\item University of Melbourne thesis on documentation cultures in Arnhem Land, 1935--2015. It covers Yolngu-led documentation generally and is not specific to sign language: \url{https://findanexpert.unimelb.edu.au/scholarlywork/1440606-documentation-cultures--arnhem-land-1935-2015}
\end{itemize}

\section{Religious and Community Resources}
No religious or community-organisation resources specific to YSL were found.

\section{Digital and Audio Resources}
\begin{itemize}
\item Photo-essay on Yolngu sign language: \url{https://phmuseum.com/projects/yolngu-sign-language?f=f}
\item Outback magazine article ``Signs of the times'': \url{https://www.outbackmag.com.au/signs-of-the-times/}
\end{itemize}

\section{Curriculum Development Context}
YSL is used mainly by adults, and the language is endangered. Any curriculum should be led by Yolngu signers and Elders, and could build on the handbook. Nothing was found on existing school programmes for deaf or hearing learners.
\end{document}